
\section{Perfect in X but not in Y}\label{sec:discussion}
Use different fonts for existence languages and operator labels:
\[
 \mathsf q=\perf_q,\quad \mathsf E=\perf_{\rm finite\ entropy},\quad
 \mathsf I=\perf(\mathrm I),
\]
\[
 \mathsf J=\perf(\II_1),\quad \mathsf U=\perf_{\sfhost},\quad
 \mathsf C=\perf_{\co}.
\]
Here $\mathsf q$ requires attainment in finite dimension; it is not the
value-one language. The inclusions are
$\mathsf q\subset\mathsf E\subset\mathsf I\subset\mathsf U\subset\mathsf C$
and $\mathsf q\subset\mathsf J\subset\mathsf U$.
Equation~\eqref{eq:languageequal} identifies $\mathsf U$ with both
$\perf(\II)$ and $\perf(\II_\infty)$; Theorem~\ref{thm:III}
identifies $\mathsf C$ with each fixed type-III language.
Write
\[
 D_2(\Sigma^0_2)=\{A\setminus B:A,B\in\Sigma^0_2\}.
\]
This is a class of differences, not the intersection of the two
pointclasses $\Sigma^0_2$ and $\Pi^0_2$. The canonical complete
language is $\{(e,f):e\in\mathrm{FIN},\ f\notin\mathrm{FIN}\}$,
where $\mathrm{FIN}$ is the $\Sigma^0_2$-complete set of indices
of finite c.e.\ sets. Indeed reduce the two $\Sigma^0_2$ predicates
separately to FIN and pair the resulting indices.

\subsection{Exact conjunction and disjunction}
\begin{backgroundlemma}[Exact game operations]
Uniform effective operations $G\wedge H$ and $G\vee H$ satisfy,
for every $\mathsf X\in\{\mathsf q,\mathsf E,\mathsf I,
\mathsf J,\mathsf U,\mathsf C\}$,
\[
 G\wedge H\in\mathsf X\iff G,H\in\mathsf X,\qquad
 G\vee H\in\mathsf X\iff G\in\mathsf X\text{ or }H\in\mathsf X.
\]
The same statements hold for each of the type languages in this paper.
\end{backgroundlemma}
\begin{proof}
First pad question and answer alphabets to common finite alphabets.
For an original zero-weight question pair accept everything; preserve
its original predicate on positive-weight pairs. Give every padded
question pair positive rational weight, accept all pairs involving
fresh questions, and reject a padded answer on an original question.
This changes no perfect-play language. It permits fixed common
question alphabets with full support; no implication below relies on
an untested pair.

For conjunction, send both players the same tag and questions for
that tagged game, and test its predicate. A perfect strategy restricts
to either tag in the same ambient algebra and state. Conversely,
tensor two perfect strategies and on a tagged question measure the
appropriate factor. Type I is closed under tensor products, as are
finite factors, type $\II_\infty$ factors, and semifinite algebras.
Amplification gives the stated $\II$ existence languages when needed.
For $\mathsf E$ use pure spatial witnesses: marginal entropies add,
so their sum is finite. Finite-dimensional witnesses stay finite.
Commuting strategies tensor as well. For $\III_\lambda$ use its
language equality with $\mathsf C$ from Theorem~\ref{thm:III}; no
closure of that subtype under direct tensoring is asserted.

For disjunction, ask the common questions and require each answer to
name a branch $s\in\{G,H\}$ as well as an answer in that branch.
Reject mismatched branches and otherwise test the named predicate.
A strategy for either game wins by always naming that branch.
For soundness let $A_{s,a}^x,B_{s,b}^y$ be the effects and put
$P_s^x=\sum_aA_{s,a}^x$, $Q_s^y=\sum_bB_{s,b}^y$.
Work first on the faithful standard-form support. The matching
argument of Section~\ref{sec:intro} makes each branch marginal a
projection, with
$P_s^x\xi=Q_s^y\xi$ for every $x,y$. Fixing $y$ and using separation
makes $P_s^x=P_s$ independent of $x$; likewise $Q_s^y=Q_s$.
These are opposite mirrors, $P_s\xi=Q_s\xi=:\xi_s$.
At least one $\xi_s$ is nonzero. Since $0\leq A_{s,a}^x\leq P_s$,
we have $A_{s,a}^xP_s=P_sA_{s,a}^x=A_{s,a}^x$, and the same holds
on Bob's side. Complete each branch POVM to the identity by adding
$I-P_s$ to one Alice answer and $I-Q_s$ to one Bob answer.
These added effects annihilate $\xi_s$. A losing branch product
annihilates $\xi$ and, by the displayed support identities, also
$\xi_s$. Normalizing $\xi_s$ therefore wins branch $s$.
All effects and the vector remain on the same Hilbert space, in the
same ambient $M,M'$; the selected vector need not be faithful.
Thus its host type is preserved. This also handles mixed strategies
by purification and standardization.

For finite entropy one can work in the original spatial purification.
Orthogonality of Bob's branch projections gives
$\rho_A=\sum_s p_s\rho_{A,s}$, $p_s=\|\xi_s\|^2$.
Entropy concavity, valid also by finite-rank truncation for infinite
entropy, gives $\sum_s p_sS(\rho_{A,s})\leq S(\rho_A)<\infty$.
Every positive-weight branch has finite entropy. This proves the last
model-specific assertion and completes both directions.
\end{proof}

\subsection{Four fixed separators and two uniform families}
Denote the companion's Weyl-clock reduction by $G^s_e$ and this
paper's selector/group reduction by $G^t_e$. Choose both reductions
for the same predicate $A(e)\equiv e\in\mathrm{FIN}$. Theorems~7.1
of \cite{Spatial} and C here, including Lemma~\ref{lem:hnn}, give
\begin{align}
 G^s_e\in\mathsf E&\iff G^s_e\in\mathsf I\iff A(e),
 &&\label{eq:sharedE}\\
 G^t_e\in\mathsf J&\iff G^t_e\in\mathsf U\iff A(e),
 &G^t_e&\in\mathsf C.\label{eq:sharedJ}
\end{align}
Every $G^s_e$ has no finite-dimensional perfect strategy; every
$G^t_e$ is commuting perfect with $\omega_q=1$.
The chained lemma additionally excludes $G^s_e$ from $\mathsf J$
and places it in $\mathsf U$ (Section~\ref{sec:sep}). The reductions
below use only its displayed spatial coordinates: $R_*$ supplies
the finite-host exclusion wherever it is needed.
No assumption about spatial perfection of a YES $G^t_e$ is made.

Take the fixed one-bridge game $R_*$ of Section~\ref{sec:shiftseparator}. Its shift model has root entropy $2$ bits and its finite compiler
factors preserve finite entropy. The finite-corner KMS contradiction
there proves that it has no finite host. Thus its complete profile
uses only the one-bridge assembly.

Let $T_*$ be the tracial solution-group game constructed in
Section~\ref{sec:sep}. It is in $\mathsf J\setminus\mathsf I$;
that section proved spatial perfection would imply finite perfection.
Let $K_{\Cat}$ be the compiled Catalan game of
\cite[Theorem C and Sections 3 and 6]{Entropy}. Its perfect spatial
strategies have infinite marginal entropy, so it is in
$\mathsf I\setminus\mathsf E$. Lemma~\ref{lem:chain} with $k=2$ applies to its nine-corner assembly:
activity $e_{g-1}+\mathsf M_g$, root mass $1/25$, transported encodings
$b_+,b_-,p,q,c$ (names of \cite[Section 6]{Entropy}), and the source-only polynomials of
\cite[Section 6]{Entropy}. Thus its extracted isometry $A$ has
modular ratio $1/4$ also on a commuting normal host. In a finite root corner $A$ would be unitary,
and KMS would give $1=1/4$; hence it is outside $\mathsf J$ as well.
Finally choose a fixed NO instance of $G^t$ and call it $Q_*$.
Theorem C gives $Q_*\in\mathsf C\setminus\mathsf U$ and
$\omega_q(Q_*)=1$. Their complete profiles are
\begin{center}\small
\begin{tabular}{@{}lcccccc@{}}
\toprule & $\mathsf q$ & $\mathsf E$ & $\mathsf I$ & $\mathsf J$ & $\mathsf U$ & $\mathsf C$\\\midrule
$R_*$ &0&1&1&0&1&1\\
$T_*$ &0&0&0&1&1&1\\
$K_{\Cat}$ &0&0&1&0&1&1\\
$Q_*$ &0&0&0&0&0&1\\\bottomrule
\end{tabular}
\end{center}
In particular $R_*\wedge T_*$ is in $\mathsf U$ but in neither
$\mathsf I$ nor $\mathsf J$, so even
$\mathsf I\cup\mathsf J\subsetneq\mathsf U$.
The remaining inclusions asserted above are strict: $R_*$ separates
$\mathsf q$ from $\mathsf E$, $K_{\Cat}$ separates $\mathsf E$
from $\mathsf I$, $T_*$ separates $\mathsf I$ from $\mathsf U$,
and $Q_*$ separates $\mathsf U$ from $\mathsf C$.
Both $\mathsf E$ and $\mathsf I$ are incomparable with $\mathsf J$.

\Needspace{16\baselineskip}
\subsection{The table, with reductions for every nonempty entry}
\begin{theorem}[Directional differences]\label{thm:differences}
The following table gives row minus column. Each suffix c means
computable many-one completeness for the displayed pointclass.
The symbol $D^\dagger$ means membership in $D_2(\Sigma^0_2)$,
$\Sigma^0_2$-hardness and $\Pi^0_2$-hardness; completeness remains open.
\begin{center}\small
\setlength{\tabcolsep}{5pt}
\begin{tabular}{@{}lcccccc@{}}
\toprule & $\mathsf q$&$\mathsf E$&$\mathsf I$&$\mathsf J$&$\mathsf U$&$\mathsf C$\\\midrule
$\mathsf q$&$\varnothing$&$\varnothing$&$\varnothing$&$\varnothing$&$\varnothing$&$\varnothing$\\
$\mathsf E$&$\Sigma^0_2$-c&$\varnothing$&$\varnothing$&$D_2$-c&$\varnothing$&$\varnothing$\\
$\mathsf I$&$\Sigma^0_2$-c&$D_2$-c&$\varnothing$&$D_2$-c&$\varnothing$&$\varnothing$\\
$\mathsf J$&$\Sigma^0_2$-c&$D_2$-c&$D_2$-c&$\varnothing$&$\varnothing$&$\varnothing$\\
$\mathsf U$&$\Sigma^0_2$-c&$D^\dagger$&$D^\dagger$&$D_2$-c&$\varnothing$&$\varnothing$\\
$\mathsf C$&$\Pi^0_1$-c&$\Pi^0_2$-c&$\Pi^0_2$-c&$\Pi^0_2$-c&$\Pi^0_2$-c&$\varnothing$\\\bottomrule
\end{tabular}
\end{center}
\end{theorem}
\begin{proof}
All empty entries follow from the stated inclusions. For upper bounds,
$\mathsf q\in\Sigma^0_1$, $\mathsf E,\mathsf I,\mathsf J,\mathsf U
\in\Sigma^0_2$, and $\mathsf C\in\Pi^0_1$.
A difference between two middle languages is therefore in $D_2$.
For a middle $\mathsf X$, $\mathsf X\setminus\mathsf q\in\Sigma^0_2$:
write its predicate as $\exists k\forall m$ for $\mathsf X$ and
conjoin the $\Pi^0_1$ predicate for nonmembership in $\mathsf q$.
Also $\mathsf C\setminus\mathsf q\in\Pi^0_1$ and
$\mathsf C\setminus\mathsf X\in\Pi^0_2$.
It remains to prove the matching lower bounds.

For the four incomparable entries use
$G^s_e\wedge R_*$ as the $\mathsf E$- or $\mathsf I$-side family,
and $G^t_f\wedge T_*$ as the $\mathsf J$-side family.
Their disjunction has $\mathsf E$ and $\mathsf I$ coordinates
exactly $A(e)$ and $\mathsf J$ coordinate exactly $A(f)$.
Thus its membership in either $\mathsf E\setminus\mathsf J$ or
$\mathsf I\setminus\mathsf J$ is $A(e)\wedge\neg A(f)$.
Swap $e,f$ between the two branches for the reverse differences.
The canonical $D_2$ problem reduces in all four cases.
This settles I-ONLY, defined here as $\mathsf I\setminus\mathsf J$.
It cannot mean absence of every type-II host, since
$\mathsf I\subset\mathsf U$.

For the $\mathsf q$ column, use $G^s_e$ for $\mathsf E$ and
$\mathsf I$; use $G^t_e\wedge T_*$ for $\mathsf J$ and
$\mathsf U$. Their larger coordinate is $A(e)$ and their
finite-dimensional coordinate is identically zero. Each such entry
is therefore $\Sigma^0_2$-hard.

For the four rightward entries of the $\mathsf C$ row, use
$G^s_e\vee Q_*$ when $\mathsf X=\mathsf E,\mathsf I$, and
$G^t_e\vee Q_*$ when $\mathsf X=\mathsf J,\mathsf U$.
The commuting coordinate is always one and the $\mathsf X$
coordinate is exactly $A(e)$, giving $\Pi^0_2$-hardness.
These instances all have $\omega_q=1$, since choosing branch $Q_*$
approximates perfection in finite dimension.
For $\mathsf C\setminus\mathsf q$, take a $\Sigma^0_1$-complete
finite-perfection reduction $M_e$ with
$M_e\in\mathsf q\iff e$ halts, available by perfect completeness
in \cite[Theorem 12.6]{MIPRE} and finite-dimensional feasibility.
Then $M_e\vee K_{\Cat}$ is always commuting perfect, has
$\omega_q=1$ because $K_{\Cat}$ is spatial perfect, and is in
$\mathsf q$ exactly when $e$ halts. Nonhalting reduces to the
difference, proving $\Pi^0_1$-hardness even under the promise.

For the two remaining complete nested differences use a shared family.
In general suppose $\mathsf Y\subset\mathsf X$ and a computable
family $\mathcal F_e$ has both coordinates equal to $A(e)$. Fix
$H\in\mathsf X\setminus\mathsf Y$. For
\begin{equation}\label{eq:sharedcriterion}
 \Lambda(e,f)=\mathcal F_e\wedge(H\vee \mathcal F_f)
\end{equation}
the $\mathsf X$ coordinate is $A(e)$ and the $\mathsf Y$
coordinate is $A(e)\wedge A(f)$. The difference coordinate is
therefore $A(e)\wedge\neg A(f)$, proving $D_2$-hardness.
Apply this first with $(\mathsf X,\mathsf Y)=(\mathsf I,\mathsf E)$,
$\mathcal F_e=G^s_e$ and $H=K_{\Cat}$, using \eqref{eq:sharedE}.
Apply it again with $(\mathsf X,\mathsf Y)=(\mathsf U,\mathsf J)$,
$\mathcal F_e=G^t_e$ and $H=R_*$, using \eqref{eq:sharedJ}.
Thus INF-ENTROPY-ONLY and $\mathsf U\setminus\mathsf J$ are
$D_2(\Sigma^0_2)$-complete. The required YES witness has finite entropy by
\cite[Theorem 7.1]{Spatial}.

Finally, for $\mathsf U\setminus\mathsf I$ and
$\mathsf U\setminus\mathsf E$, the map $e\mapsto G^t_e\wedge T_*$
is $\Sigma^0_2$-hard: its $\mathsf U$ coordinate is $A(e)$ and
both spatial coordinates vanish. The map $e\mapsto G^s_e\vee T_*$
is $\Pi^0_2$-hard: its $\mathsf U$ coordinate is one and both
$\mathsf E,\mathsf I$ coordinates equal $A(e)$.
Together with the $D_2$ upper bound, this proves exactly the two
$D^\dagger$ entries. The theorem makes no completeness claim for them.
\end{proof}

\begin{corollary}\label{cor:entropyonly}
The language of games with perfect spatial play but no finite-entropy
perfect spatial strategy is $D_2(\Sigma^0_2)$-complete.
In particular it belongs to neither $\Sigma^0_2$ nor $\Pi^0_2$.
\end{corollary}
\begin{proof}
The identification is $\mathsf I\setminus\mathsf E$ in the theorem.
A $D_2$-complete language is hard for both $\Sigma^0_2$ and
$\Pi^0_2$ by fixing one of the canonical FIN indices to a YES or NO
index. Membership in either pointclass would collapse these two
pointclasses, contradicting strictness of the arithmetical hierarchy.
\end{proof}
The corollary concerns exact attainment; it asserts no rate at
which entropy must diverge for near-optimal strategies.

For differences between middle models, value-one padding can
be made explicit. Replacing $G$ by $G\vee Q_*$ preserves all lower
coordinates and sets the commuting coordinate to one, while forcing
$\omega_q=1$. To remove finite-dimensional perfect play as well,
use $G\wedge(R_*\vee T_*)$; the second factor has all five
infinite-model coordinates equal to one and its $\mathsf q$
coordinate zero. In particular \eqref{eq:sharedcriterion} for
INF-ENTROPY-ONLY already has no finite-dimensional perfect instance;
adding $Q_*$ also makes every instance value one.
The $\mathsf C$-row hardness families above already have value one,
including the separate $\mathsf C\setminus\mathsf q$ construction.

\subsection{The remaining nested differences}
For a nested pair $\mathsf Y\subset\mathsf X$ of $\Sigma^0_2$-complete
languages with a separator $H$, independent hardness families give
both kinds of hardness: $G_X(e)\wedge H$ has difference predicate
$A(e)$ and $G_Y(e)\vee H$ has difference predicate $\neg A(e)$.
They do not by themselves provide the shared family required by
\eqref{eq:sharedcriterion}. This explains the two open table cells.
We use II-ONLY to mean $\mathsf U\setminus\mathsf I$, with
$\mathsf U\setminus\mathsf E$ its finite-entropy variant; the
incomparable difference $\mathsf J\setminus\mathsf I$ is already
complete by the theorem.

\begin{remark}[A finite monotone obstruction]
An exhaustive Boolean check in the verification bundle ranges over all 7581
monotone Boolean functions of five inputs
$G^t_e,G^t_f,G^s_e,G^s_f,T_*$ and all profiles consistent with the
proved coordinates. For $\mathsf U\setminus\mathsf I$, none realizes
$A(e)\wedge\neg A(f)$ uniformly over the unproved spatial coordinates
of $G^t$. Even setting both of those coordinates optimistically to
zero gives none. With a shared family $\mathcal F_e,\mathcal F_f$ as two inputs,
351 functions do realize it. This is a bounded exhaustive check,
not an impossibility theorem for new game constructions. Separate completeness and the inclusion lattice do not supply
the shared-family equivalence needed by this reduction.
\end{remark}

An open route is a selector-fed Fock tableau, replacing the Weyl
clock by a 16-ary critical selector. One would require commuting
selector and tableau labels and an equality reader
\[
 (2C_{\rm ctr}-C_{\rm sel})P_{\rm ret}=0,\qquad
 T_{\rm inv}(2C_{\rm ctr}-C_{\rm sel})
 =\tfrac18 C_{\rm ctr}C_{\rm sel}(I-P_{\rm ret}),
 \qquad P_HP=0.
\]
The intended pulse selects a finite candidate, and a finite-depth
selector tensored with the Fock tableau would be the finite-entropy
YES model. What remains open is soundness of the additional selector/tableau
relations with the extra defect-isometry family, and the resulting
shared-family equivalence
$\mathsf E(\mathcal F_e)=\mathsf I(\mathcal F_e)=\mathsf U(\mathcal F_e)=A(e)$.
These equations describe the proposed shared-family construction. Completeness for both
II-ONLY variants therefore remains open.

\section{Scope and remaining questions}\label{sec:scope}
These are existence languages: the host and its normal state may be
chosen together. They do not decide whether a prescribed factor and
state realize a correlation, or determine the least resources of a
strategy. Beyond the two open nested differences, natural questions
ask which syntactic classes force finite-entropy spatial perfection
whenever spatial perfection is possible, and how quantitative
entropy requirements interact with the difference classification.
The finite-algebra YES model of the present selector/group family
alone gives no spatial YES model; replacing that missing assertion
by separate completeness would not settle either open cell.
